

\noindent \textcolor{azulp}En clase se describió el Modelo Geométrico Browniano para modelar el precio de los activos:
\[
dS_t = \mu S_t dt + \sigma S_t dW_t
\]

\noindent Explique con palabras \textbf{\textit{no técnicas}} los siguientes componentes de dicha ecuación:

\vspace{0.3cm}
\begin{itemize}[label=\textcolor{azulp}{$\circ$}, itemsep=0.4cm, leftmargin=1.5cm]
    \item $dS_t$: \textbf{El cambio inmediato en el precio.} Representa cuánto sube o baja el valor del activo en un instante de tiempo muy, muy pequeño. Es la "foto" del movimiento del precio de un momento a otro.
    
    \item $\mu S_t dt$: \textbf{La tendencia esperada (lo predecible).} Representa hacia dónde se dirige el precio de forma "natural" con el paso del tiempo, asumiendo que el mercado fuera completamente tranquilo. Es como la "velocidad crucero" o el rendimiento promedio.
    
    \item $\sigma S_t dW_t$: \textbf{El factor sorpresa (lo impredecible).} Representa los choques aleatorios, las noticias y el "ruido" del mercado que hace que el precio salte hacia arriba o hacia abajo de forma inesperada. Es el riesgo o la volatilidad pura.
    
    \item $\mu S_t dt + \sigma S_t dW_t$: \textbf{La realidad completa del mercado.} Nos dice que el movimiento total del precio de un activo es simplemente la suma de dos cosas: su tendencia esperada (lo que dicta la lógica) más la volatilidad del mercado (el caos del día a día).
\end{itemize}
